
Jusqu'à présent nous avons utilisé l'ordinateur pour
piloter le \robot\ en mode {\em interactif}\index{Mode interactif} ou
{\em conversationnel}\index{Mode conversationnel}. Dans ce mode de fonctionnement 
nous tapons sur une touche du clavier de l'ordinateur, qui transmet
une commande au \robot\ que celui-ci exécute immédiatement à moins que par erreur
nous tentions de le faire sortir du cadre auquel cas il répond par
un \raisebox{0.25ex}{$\scriptscriptstyle \boldmath\ll$}~%
bip-bip%
~\raisebox{0.25ex}{$\scriptscriptstyle \boldmath\gg$}. 
Nous recommençons ainsi, en corrigeant les fautes de pilotage au fur et à mesure que
le dessin se construit, jusqu'à ce que nous
indiquions à l'ordinateur que le dessin est terminé. 

\smallskip

Alors que la commande interactive du \robot\ correspond à l'exécution
im\-mé\-dia\-te d'une tâche, la commande programmée du \robot\ consiste
à séparer la description de la tâche de ses exécutions ultérieures. 

\smallskip

Plus
généralement, l'idée de la {\em programmation}\index{Programmation} est la suivante~: 
pour faire quelque chose, le {\em programmeur} écrit un {\em programme}\index{Programme}
(dans un certain {\em langage}\index{Langage de programmation}) qui décrit avec précision comment 
effectuer cette chose. Ensuite il demande à un {\em ordinateur}\index{Ordinateur} d'exécuter
le programme, en suivant exactement les instructions données par le
programme. Autrement dit, au lieu de faire la chose
lui-même tout de suite, le programmeur explique à un ordinateur comment 
le faire à sa place plus tard$\:$! On évite ainsi d'avoir à réaliser
des tâches longues et répétitives, à condition de savoir 
décrire ces tâches con\-ve\-na\-ble\-ment.

\smallskip

S'il s'est trompé, le programmeur
doit corriger ses fautes en modifiant le programme. C'est ce qu'on appelle 
{\em déboguer}\index{Déboguer} le programme. Le {\em déboguage}\index{Déboguage}
du programme consiste à trouver puis à éliminer les {\em bogues}\index{Bogue} c'est-à-dire les 
erreurs se manisfestant à l'exécution par des anomalies de fonctionnement (comme par 
exemple {\idt pqtg} au lieu de {\idt pqtd} ou {\idt st} oublié à la fin du programme)%
\footnote{\ Nous utilisons le vocabulaire recommandé par la commission
ministérielle de ter\-mi\-no\-lo\-gie \cite{terminologie-info}. La bogue est
la traduction homonymique en français du mot américain {\em bug}\index{Bug} qui 
désigne les petits insectes nuisibles et plus familièrement les microbes 
et autres virus aussi bien que les petites choses qui clochent comme les erreurs 
dans les programmes. La bogue désignant en français l'enveloppe piquante
de la châtaigne, l'idée de désagrément est conservée. L'enveloppe fait
également penser à la coquille typographique. La traduction
ne garde pas l'idée de petit grain de sable aux conséquences dramatiques
du mot américain.}.
Le principe de fonctionnement des ordinateurs commandé par un programme a des 
analogies avec de nombreuses activités humaines
comme le montre la table \hbox{ci-dessous~:}

\begin{table}[htbp]
\begin{center}
\small
\begin{tabular}{@{}|p{25mm}|p{30mm}|p{30mm}|p{30mm}|@{}}
\hline
{\em Programmeur } & {\em Programme } & {\em Ordinateur } & {\em Dessin }  \\ 
\hline \hline
Chef-cui\-si\-nier & Re\-cet\-te de cui\-si\-ne & Cui\-si\-niers & Re\-pas\\
Com\-po\-si\-teur & Œu\-vre mu\-si\-ca\-le & Mu\-si\-ciens & Con\-cert \\
Au\-teur & Tex\-te d'une piè\-ce de thé\-â\-tre  & Ac\-teurs & Re\-pré\-sen\-ta\-tion \hfill thé\-â\-tra\-le \\
Cho\-ré\-gra\-phe & Cho\-ré\-gra\-phie & Dan\-seurs & Bal\-let \\
Réa\-li\-sa\-teur & Film & Pro\-jec\-teur \hfill ci\-né\-ma\-to\-gra\-phi\-que & Sé\-an\-ce de ci\-né\-ma \\
Tis\-se\-rand & Car\-tes per\-fo\-rées & Mé\-tier \hfill à \hfill tis\-ser jac\-quard & Tis\-su jac\-quard\\
%Ma\^\i\-tre-Êcu\-yer & Tex\-te de re\-pri\-se & Ca\-va\-lier et che\-val & Re\-pri\-se \hfill de dres\-sa\-ge \\
Chef de pis\-te & Plan du par\-cours & Ca\-va\-lier et che\-val & Par\-cours de saut d'obs\-ta\-cles \\
\hline
\end{tabular}
\end{center}
%\label{tab:act-prog}
\caption{Exemples d'activités programmées}
\end{table}

Nous allons maintenant apprendre à programmer les ordinateurs
en \pascal. Le pilotage du \robot\ pour faire des dessins sera notre 
principale application. Pour faire nos dessins dans
ce mode de fonctionnement programmé, nous procédons comme suit~:
\begin{enumerate}
\item Nous commençons par écrire un programme \pascal\ où nous 
in\-di\-quons quelle est la séquence de commandes que le \robot\ devra 
exécuter pour réaliser le dessin que nous avons en tête ou que
nous avons pré\-a\-la\-ble\-ment dessiné sur une feuille de papier
quadrillé. Ce programme est écrit avec l'aide de l'ordinateur
(avec un éditeur de texte) qui le conserve sur disque ou sur disquette.
\item  Ensuite nous demandons à l'ordinateur de {\em compiler}\index{Compiler} ce 
programme \pascal, c'est-à-dire de le traduire dans son 
{\em langage machine}\index{Langage machine}.
Le programme ainsi traduit s'appelle le {\em programme compilé}\index{Programme compilé}. 
Le programme compilé et le programme
\pascal\ ori\-gi\-nal sont équivalents en ce sens qu'ils décrivent 
exactement la même séquence de commandes pour le \robot. Ils sont
simplement écrits dans des langages différents.
\item  Enfin nous demandons à l'ordinateur  d'{\em exécuter}\index{Exécution} le 
programme compilé. Ce faisant, l'ordinateur pilote le \robot\ en lui
demandant d'exé\-cu\-ter les commandes dans l'ordre où nous les avons 
inscrites dans le programme \pascal\ ori\-gi\-nal. 
\end{enumerate}

{ %1
\newcommand{\touche}[1]{{
\begin{picture}(12,12)(0, 2)
\put(0,0){\framebox(12,12)[lb]{}}
\put(1.5,2){\makebox(0,0)[lb]{\tiny #1}}
\end{picture}\hspace{0.3ex}}}
\begin{exemple}[Carré]{\rm
Soit à dessiner le carré ci-contre. Le programme
\pascal\ sera le suivant~:

{\raggedright 
 \vspace{\baselineskip}
\vspace*{1.3cm}%
\includevspacePostScript[-\baselineskip]{PetitCarre}{0.95}{0.95}%
\vspace*{-1.3cm}%
 \input Carre.tex
 }
\noindent Quand il exécute ce programme après l'avoir compilé, 
l'ordinateur transmet au \robot\ l'ordre d'exécuter la séquence
 de commandes {\idt av} puis {\idt av} puis
{\idt pqtd} puis {\idt av} \ldots\ et enfin {\idt pqtd} exactement comme si 
nous avions tapé au clavier \touche{A} puis successivement \touche{A},
\touche{+}, \touche{A}, \touche{A}, \touche{+}, \touche{A}, \touche{A},
 \touche{+}, \touche{A}, \touche{A} et enfin \hbox{\touche{+}\@.} La dernière commande
 {\idt st} indique à l'ordinateur qu'il doit attendre que nous tapions
 sur la touche de retour à la ligne 
 \framebox[11.mm][l]{$\hookleftarrow$~~~} pour terminer l'exécution du
 programme. Ceci nous laisse le temps de contempler à loisir le dessin
 du \robot\ avant qu'il ne disparaisse. Le dessin n'est pas perdu puisqu'il
 suffit d'exécuter à nouveau le programme pour le reproduire.
 \noindent $\Box$ }\end{exemple}
 } %1

On constate sur ce premier exemple que le langage \pascal, bien que
conçu par le professeur suisse allemand Niklaus Wirth de la 
Eidgenössische Technische Hochschule à Zürich, utilise des mots
anglais comme `{\mcl program}' (qui signifie `pro\-gram\-me'), `{\mcl uses}'
 (qui
signifie `uti\-li\-se' et se prononce ``iou\-seu\-zzz'' avec l'accent anglais 
évidemment),  `{\mcl be\-gin}' (qui signifie `début' et se prononce très 
approximativement ``biguinne'') et `{\mcl end}' (qui signifie `fin' et qui,
Outre-Manche, se prononce ``ainde''). L'informatique conserve les traces de son origine 
anglo-\hbox{saxonne$\:$!}

Pour l'instant il nous suffit de savoir
que l'ordinateur qui exécute ce programme demande au \robot\ d'exécuter 
les commandes figurant entre le {\mcl begin} et le {\mcl end}, les unes après
les autres. Les
points-virgules qui séparent ces commandes peuvent être compris comme signifiant `puis'.
  
 
%%%%%%%%%%%%%%%%%%%%%%%%%%%%%%%%%%%%%%%%%%%%%%%%%%%%%%%%%%%%%%%%%%%%%%%%%%%
\noindent \rule{137mm}{0.1mm} %---------------------------------------------
\begin{exercice}{
%\label{exo:trouver-cocotte}
Exécuter le programme {\idt Des\-si\-ner} de commande 
interactive du \robot\ et lui donner les ordres décrits par 
le programme \pascal\ suivant (la correspondance entre les touches du
clavier et les commandes du \robot\ est donnée par la table
\ref{tab:cmd-inter} page \pageref{tab:cmd-inter})~:
\input Inconnu.tex
\noindent Quel dessin avez-vous demandé au \robot\ de réaliser$\:$?
\noindent $\Box$}\end{exercice}
\begin{exercice}{
%\label{exo:corriger-E}
Nous voudrions que le programme ci-dessous dessine la lettre {\tt E} majuscule 
\hbox{ci-contre~:}

{\raggedright 
\vspace{\baselineskip}
\vspace*{3cm}%
\includevspacePostScript[-\baselineskip]{E-epaisse}{0.95}{0.95}
\vspace*{-3cm}%
 \label{prog:LettreEavecBogues}
 \input LettreEavecBogues.tex
 }
\noindent Ce programme erroné contient deux bogues. Pouvez-vous les trouver$\:$?
\noindent $\Box$}\end{exercice}
\begin{exercice}{
%\label{exo:exec-stegosaure}
Exécuter le programme \pascal\ {\idt Stegosaure} figurant sur la disquette
d'accompagnement de ce livre.
\noindent $\Box$}\end{exercice}
%%%%%%%%%%%%%%%%%%%%%%%%%%%%%%%%%%%%%%%%%%%%%%%%%%%%%%%%%%%%%%%%%%%%%%%%%%%
\begin{corrige}{
%\label{cor:trouver-cocotte} 
Une cocotte en papier.
\noindent $\Box$}\end{corrige}
\begin{corrige}{
%\label{cor:corriger-E} 
Le programme {\idt LettreEavecBogues} de la page \pageref{prog:LettreEavecBogues} 
contient deux bogues~:\\
\begin{minipage}[t]{115mm}
\begin{enumerate}
\item la troi\-si\-ème commande {\idt pqtd} de la deuxième ligne
devrait être une commande {\idt pqtg},
\item il manque une commande {\idt av} à l'avant-dernière ligne. 
\end{enumerate}
Le résultat de l'exécution de ce programme est reproduit ci-contre
(le \robot\ étant placé dans sa position finale).
Le programme débogué qui dessine correctement la lettre {\tt E} épaisse est le suivant (les corrections sont 
sou\-li\-gnées)~:
\end{minipage}
\hfill
\begin{minipage}[t]{17mm}
\vspace{3.2cm}\includePostScript{lettre-e-avec-bogues}{0.95}{0.95}%
\hspace{-1.7cm}  % 1.7cm de large
\end{minipage}
\input LettreE.tex
\noindent $\Box$}\end{corrige}
\begin{corrige}{
%\label{cor:exec-stegosaure}
La façon d'exécuter le programme \pascal\ {\idt Ste\-go\-sau\-re}
dépend de l'ordinateur \hbox{utilisé~:}

\tiret Sur le 
Macintosh, il faut cliquer deux fois rapidement sur l'icône {\idt Ste\-go\-sau\-re.p}
puis choisir l'option `{\tt Exécuter}' du  menu `{\tt Compil.}'. 

\tiret Sur les
compatibles IBM PC, le programme {\idt Stegosaure} étant rangé dans le
fichier {\idt Stegosaure.pas}, taper {\idt tpc Stegosaure} pour compiler le
programme \pascal\ en un programme compilé {\idt Stegosaure.exe} que l'on
exécute en tapant {\idt Stegosaure}.

\noindent $\Box$}\end{corrige}

%%%%%%%%%%%%%%%%%%%%%%%%%%%%%%%%%%%%%%%%%%%%%%%%%%%%%%%%%%%%%%%%%%%%%%%%%%%
\noindent \framebox[\textwidth]{
\begin{minipage}{\myparboxwidth}
\setlength{\parindent}{2em}
\small \sf
\vspace{\abovedisplayskip}
\section*{Addenda}
Nous devons à Seymour Papert, l'auteur de {\sc Logo} \cite{logo}, la diffusion de l'idée 
d'utiliser le dessin pour apprendre les concepts de base de la programmation. Cependant
l'introduction d'un quadrillage nous évite l'usage des angles et des longueurs comptées en 
pixels\@%
\footnote{\ Le {\em pixel}\index{Pixel} est le plus petit
élément de l'écran graphique qui peut être allumé (blanc), éteint (noir) ou
colorié avec des luminosités et des couleurs plus ou moins nombreuses selon
la qualité de l'écran.}. Le quadrillage permet de tracer très facilement
des arcs de
cercles, ce qui n'est pas immédiat en {\sc Logo}.

L'idée d'utiliser un 
quadrillage pour faire des dessins géométriques n'est pas nouvelle. On la connaît à 
l'école depuis longtemps (voir \cite{thevenet1} par exemple). On la retrouve,
à la base d'exercices de programmation graphique, dans
de nombreux ouvrages informatiques de haut niveau comme en \TeX\@\footnote{\ \TeX\ ou plutôt sa variante \LaTeX\ 
\cite{lamport1} est le langage de programmation que nous avons utilisé pour
composer ce livre.}\ par exemple (voir \cite{knuth1} pages 389--390). 
Le quadrillage
élimine les problèmes de représentation
discrète sur un écran fini de concepts mathématiques continus. Par conséquent
nos segments ressemblent à des segments tels qu'on les dessine sur le papier 
(et non pas à des traits juxtaposés en escalier), nos
cercles ressemblent à des cercles (et non pas
à des carrés comme cela arrive sur les écrans graphiques pour des cercles trop
petits).  Certes l'emploi de grands carreaux limite les possibilités 
de dessin. La programmation n'en est que plus facile si nous tirons avantage
des formes géométriques pour induire la structure des programmes.
Une fois que les notions de base seront acquises nous pourrons faire varier la taille de la grille
de sorte qu'en passant à la limite fixée par les pixels de l'écran
nous retrouverons toutes les possibilités offertes par les écrans graphiques actuels 
à haute résolution. 

Dans le chapitre suivant nous commençons à écrire nos premiers
programmes \pascal. Les algorithmes\@\footnote{\ Un {\em algorithme} est la
description d'un enchaînement d'opérations sur des ob\-jets pour 
accomplir une tâche. Ce mot a pour origine la version latine du nom du 
mathématicien et astronome persan al-Khwarizmi du IX\raisebox{0.7ex}{\tiny ème} 
siècle de notre ère.
Son traité {\em Al-Jabr wa l-muqabala} enseigne comment résoudre algorithmiquement
les équations du premier et second degré à cœfficients numériques.
Le terme d'{\em algèbre\/} dérive du titre même de ce traité.
} restent très simples puisqu'ils
consistent à codifier le mouvement de la main qui dessine.

\vspace{\abovedisplayskip}
\end{minipage}
}
%%%%%%%%%%%%%%%%%%%%%%%%%%%%%%%%%%%%%%%%%%%%%%%%%%%%%%%%%%%%%%%%%%%%%%%%%%%
